\documentclass[11pt,a4paper]{article}
\usepackage[margin=25mm]{geometry}
\usepackage{amsmath,amssymb,amsthm,hyperref}
\setlength{\emergencystretch}{2em}
\newtheorem{theorem}{Theorem}
\newcommand{\gen}[1]{\langle#1\rangle}
\title{Lifting an odd singleton budget across the opposite end}
\author{Human Compression project: L044 continuation}
\date{}
\begin{document}\maketitle
Let $G=\gen{\rho_0,\ldots,\rho_{r-1}}\le A_n$ be a transitive
connected string C-group, $n=2r\ge14$, with every maximal omission
intransitive. Suppose, as actual permutation subgroups, that
\[
 G_0=\operatorname{Alt}(\Omega\setminus\{a\}),\qquad
 C=G_{r-1}=\operatorname{Alt}(\Omega\setminus\{d\}).       \tag{1}
\]
Assume every colour $2,\ldots,r-2$ has at least two ambient-eligible
edges. An ambient-eligible $j$-edge joins distinct $G_j$-orbits.
All the previously accepted ambient interval and selected-localization
proofs, including their complete retained inputs, are required.

\begin{theorem}
Under these hypotheses either $G_{>1}$ fixes at least three points
or $G_{>2}$ fixes at least five. Equivalently there is $h\in\{1,2\}$
such that, for $X_h=\operatorname{supp}(G_{>h})$,
\[
                         2h\le n-|X_h|-1.              \tag{2}
\]
The fixed points in this conclusion belong to the full ambient tail,
not merely the corresponding subgroup of $C$.
\end{theorem}
\begin{proof}
Put $\eta=\rho_{r-1}$, $\sigma=\rho_1$, $b=a^{\rho_0}$,
\[
 H_C=G_{\{2,\ldots,r-2\}},\quad
 K_C=G_{\{3,\ldots,r-2\}},\quad
 H=\gen{H_C,\eta},\quad K=\gen{K_C,\eta}.
\]
The points $a,d$ are distinct, since otherwise $G$ fixes their common
point. The generator $\eta$ moves $d$ into the natural $C$-orbit,
and fixes $a$. Because it commutes with $\rho_0$, it also fixes $b$.
Similarly $\rho_0$ fixes $d$. Thus $\rho_0=(a\ b)\delta$, where
$\delta$ is an odd involution fixing $a,b,d$, and it centralizes
$H_C$ and $\eta$. The unique ambient-eligible $\eta$-edge is the
one incident with $d$.

We first justify the use of the accepted odd singleton theorem in
$C$. Its degree is $n-1=2(r-1)+1$. Its marking is connected:
otherwise IP decomposes this nonabelian simple group as a direct
product of nontrivial marked factors. For any omitted colour
$j\le r-2$, if $C_j$ were transitive on $\Omega\setminus\{d\}$,
then $\gen{C_j,\eta}\le G_j$ would be transitive on $\Omega$, which
is excluded. Its first omission is the actual natural alternating
group on $\Omega\setminus\{a,d\}$, by IP and (1).
Every specified ambient-eligible edge lies in the $C$-domain and
is also $C_j$-eligible, since $C_j\le G_j$. This is only the valid
global-to-local direction.

Let $f_C$ be the number of $H_C$-fixed points in the $C$-domain,
and let $c$ be its nontrivial orbit count. The local interval bound
and the accepted odd theorem give
\[
 f_C+c\le5,\qquad f_C\ge3\quad\hbox{or}\quad
 \left\{\begin{array}{l}
 f_C=2,\ H_C\hbox{ has orbits }\{a\},\{b\},E,V,\\
 E=\{u,v\},\ |V|=n-5\hbox{ odd},\
 \rho_0=(a\ b)(u\ v).
 \end{array}\right.                                   \tag{3}
\]
Here $c\ge1$, since the marked subgroup $H_C$ is nontrivial.
We now analyze the extra fixed points rather than assume they
remain fixed on adjoining $\eta$.

Suppose $f_C=3$, with additional fixed point $z$.
The involution $\delta$ fixes $z$. If it is free on the entire
nontrivial $H_C$-support (including the possibility of exchanging
two orbits), its full fixed set is $\{a,b,z,d\}$. Since $\eta$
centralizes $\delta$, fixes $a,b$ and moves $d$, it swaps $d,z$.
But $\sigma$ fixes $d$ and centralizes $\eta$, so fixes $z$ too.
Consequently $G_0=\gen{\sigma,H_C,\eta}$ preserves $\{d,z\}$,
contrary to (1).

The only other possibility has two preserved nontrivial orbits:
$O$ of size $2q$, on which $\delta$ is free, and $V$ of even
size $v$, on which it is the identity. Oddness makes $q$ odd.
Partition the colours $2,\ldots,r-1$ into $L$, those with an
ambient-eligible edge in $O$, and $F$, the rest. The projected
$\delta$-pair forest gives $|L|\le q-1$: any alternate path of
pair-blocks lifts using $\rho_0\in G_j$ and violates the selected
ambient omission cut. The colour $r-1$ belongs to $F$, since its
only eligible edge is incident with the $\delta$-fixed point $d$.
Every eligible edge of $F\setminus\{r-1\}$ is confined to
$W=V\cup\{z,d\}$, which has three $H_C$-orbits, including both
singletons. Selected localization in the full interval
$\{2,\ldots,r-2\}$ gives
\[
 2(|F|-1)\le |W|-o_{H_C}(W)=v-1,
 \qquad |F|\le v/2.
\]
Thus $r-2=|L|+|F|\le q-1+v/2$, contradicting
$n=4+2q+v=2r$. This excludes $f_C=3$.

If $f_C=4$, then $c=1$ and its nontrivial orbit has odd size.
The centralizing involution $\delta$ fixes that orbit pointwise.
Its oddness therefore forces it to swap the other two fixed points,
say $E=\{u,v\}$. Again $\rho_0=(a\ b)(u\ v)$, but now $H_C$
fixes $E$ pointwise. If $\eta$ also fixes $E$, then $H$ has at
least four fixed points and (2) holds with $h=1$.
Values $f_C\ge5$ are impossible by (3) and $c\ge1$.

It remains to consider the literal double-transposition head from
either surviving case. In both cases $H_C$ preserves $E$ and is
transitive on the odd remaining set $V$. In the $f_C=2$ case,
the second proof of the accepted odd theorem proves explicitly
that $K_C$ fixes $E$ pointwise; if it did not, the two pair images
and the conjugate $\rho_0^\sigma$ would confine all its eligible
edges to a set of size exactly twice its rank, contradicting strict
localized counting. In the $f_C=4$ case this fixation is immediate.
Hence in both cases $K_C$ fixes
\[
                         a,b,E,E^\sigma,\ x=b^\sigma.   \tag{4}
\]
Transitivity of $C_0=\gen{\sigma,H_C}$ gives $x\ne b$;
$\sigma$ fixes $a,d$, so $x\in E\cup V$.
Also $\eta$ fixes $x$, since it centralizes $\sigma$ and fixes $b$.

Suppose $\eta$ fixes $E$ pointwise. It then fixes $E^\sigma$.
If $x\in V$, (4) gives five distinct $K$-fixed points
$a,b,u,v,x$. If $x\in E$, rename it $u$; then $u^\sigma=b$,
and $w=v^\sigma$ belongs to $V$. Indeed $w=v$ would make
$\{b,u,v\}$ invariant under $C_0$, and the other excluded values
follow from injectivity and fixation of $a,d$.
Now $a,b,u,v,w$ are five distinct $K$-fixed points.
This gives (2) with $h=2$.

Finally suppose $\eta$ swaps $E$. Then $x\notin E$.
The pair $E^\sigma$ is also $\eta$-invariant. It cannot equal $E$,
as that would make $E$ invariant under $C_0$. Their intersection
cannot be a singleton, because $\eta$ would fix that singleton
and hence both points of $E$. Thus they are disjoint, and
$E^\sigma=\{w,z\}\subset V\setminus\{x\}$.
By (4), $K_C$ fixes the seven distinct points
$a,b,x,u,v,w,z$ in the $C$-domain, as well as $d$.
All its actual edges, and hence all its ambient-eligible edges,
are confined to the invariant complement $Y$ of these eight points.
Every one of its $r-4$ colours has two such eligible edges.
But $|Y|=n-8=2(r-4)$, and the localized interval theorem gives
\[
              2(r-4)\le |Y|-o_{K_C}(Y)<|Y|,
\]
a contradiction. This completes all cases and proves the ambient
fixed-support budget.
\end{proof}

\paragraph{Attribution and remaining scope.}
This supplies a precise two-singleton-end embedding step for
Cameron--Fernandes--Leemans--Mixer, \emph{Highest rank of a polytope
for $A_n$}, arXiv:1605.09173v1, Proposition 5.32 (physical pages
44--45). The source's fixed-support budget strategy and the complete
accepted odd singleton and selected-localization proofs are reused.
The additional argument proves which local fixed points survive the
opposite generator; it makes no general parabolic inheritance claim.
The two natural alternating omissions, the even critical degree and
the paired interior colours are explicit hypotheses. Other opposite
heads, final budget combination, source equality diagrams and the
remaining rigidity/table contracts are not recovered here. No new
computation, formal statement, original or question credit is claimed.
\end{document}
